\subsection{局部紧交换群上的正定函数}

在本节中，G 是局部紧交换群，\(\mu\) 表示 G 上的哈尔测度。以 \(\widehat{G}\) 表示 G 的对偶群，以 \(\widehat{\mu}\) 表示 \(\mu\) 的对偶哈尔测度（II, p. 214, 定义 4）。以 F 表示 G 上有界测度的巴拿赫空间 \(\mathcal{M}^{1}(G)\) 上的傅里叶变换。

由于 \(\mathcal{C}_0(\mathrm{G})\) 是 \(\mathcal{K}(\mathrm{G})\) 在 \(\mathcal{C}_b(\mathrm{G})\) 中的闭包，巴拿赫空间 \(\mathcal{M}^1 (\mathrm{G})\)（\(\mathcal{K}(\mathrm{G})\) 赋予 \(\mathcal{C}_b(\mathrm{G})\) 的拓扑的对偶）与 \(\mathcal{C}_0(\mathrm{G})\) 的对偶相等同（参见 INT, III, p. 56, § 1, no. 8）。在 \(\mathcal{M}^1 (\mathrm{G})\) 上赋予与此对偶性相联系的弱拓扑。

引理 9. --- 傅里叶变换是从赋予弱拓扑的 \(\mathcal{M}^{1}(\mathrm{G})\) 到 \(\mathcal{C}_{b}(\widehat{\mathrm{G}})\) 的连续映射，后者赋予由弱拓扑 \(\sigma(\mathrm{L}^{\infty}(\widehat{\mathrm{G}}),\mathrm{L}^{1}(\widehat{\mathrm{G}}))\) 诱导的拓扑。

设 \(\mathfrak{F}\) 是 \(\mathcal{M}^1 (\mathrm{G})\) 上弱收敛于 G 上有界测度 \(\nu\) 的滤子。对每个 \(\varphi \in \mathrm{L}^1 (\widehat{\mathrm{G}})\)，有 \(\mathcal{F}(\varphi)\in \mathcal{C}_0(\mathrm{G})\)（II, p. 209, 命题 4），从而

\[
\begin{array}{r l} & {\int_ {\widehat {\mathrm{G}}} \varphi \mathcal {F} (\nu) d \widehat {\mu} = \int_ {\mathrm{G}} \mathcal {F} (\varphi) d \nu} \\ & {\qquad = \lim _ {\eta , \mathfrak {F}} \int_ {\mathrm{G}} \mathcal {F} (\varphi) d \eta = \lim _ {\eta , \mathfrak {F}} \int_ {\widehat {\mathrm{G}}} \varphi \mathcal {F} (\eta) d \widehat {\mu},} \end{array}
\]

由傅里叶变换的转置性质（II, p. 221, 命题 13）。引理得证。

定理 5（博赫纳）。--- 连续函数 \(\varphi\) 在 \(\widehat{\mathbf{G}}\) 上且属于 \(\operatorname{Pos}(\widehat{\mathbf{G}})\)，当且仅当存在 \(\nu\) 在 \(\mathbf{G}\) 上的有界正测度，使得 \(\varphi = \mathcal{F}(\nu)\)。

设 \(\nu \in \mathcal{M}^{1}(\mathrm{G})\) 是正测度。其傅里叶变换连续。对每个有限族 \((x_{i})_{i\in \mathrm{I}}\)（属于 \(\widehat{\mathbf{G}}\)）以及复数的每个有限族 \((t_i)_{i\in \mathrm{I}}\)，可得

\[
\begin{array}{r} \sum_ {i \in \mathrm{I}} \sum_ {j \in \mathrm{I}} \bar {t} _ {i} t _ {j} \mathcal {F} (\nu) (x _ {i} ^ {- 1} x _ {j}) = \sum_ {i \in \mathrm{I}} \sum_ {j \in \mathrm{I}} \bar {t} _ {i} t _ {j} \int_ {\mathrm{G}} x _ {i} \overline {{x}} _ {j} d \nu \\ = \int_ {\mathrm{G}} \left| \sum_ {i \in \mathrm{I}} \bar {t} _ {i} x _ {i} \right| ^ {2} d \nu \geqslant 0, \end{array}
\]

由于 \(\nu\) 是正测度。因此，\(\nu\) 的傅里叶变换是 \(\hat{G}\) 上的正定函数（V, p. 432, 定理 1, (ii)）。

证明逆向蕴含。给集合 \(\operatorname{Pos}_{\leqslant1}(\widehat{\mathbf{G}})\) 赋予弱拓扑，该拓扑像前面一样由弱拓扑 \(\sigma(\mathrm{L}^{\infty}(\widehat{\mathrm{G}}),\mathrm{L}^{1}(\widehat{\mathrm{G}}))\) 诱导。集合 \(\operatorname{Pos}_{\leqslant1}(\widehat{\mathbf{G}})\) 是紧凸集（V, p. 448, 命题 13 的推论）。由 V, p. 450, 命题 14，V, p. 444, 命题 11 以及 V, p. 390, 推论 7，\(\operatorname{Pos}_{\leqslant1}(\widehat{\mathbf{G}})\) 的极点是零函数和 \(\widehat{G}\) 中的元素。

令 N 为 G 上质量 \(\leqslant 1\) 的正有界测度的集合；它在赋予弱拓扑的 \(\mathcal{M}^{1}(G)\) 中是紧的（EVT, III, p. 17, 推论 3）。由引理 9 及证明的第一部分，G 上的傅里叶变换通过对子空间的传递定义了从 \(\mathcal{N}\) 到 \(\mathrm{Pos}_{\leqslant 1}(\widehat{\mathrm{G}})\) 的连续映射；由齐次性，只需证明此映射是满射。\(\mathcal{F}(\mathcal{N})\) 是傅里叶变换下 \(\mathcal{N}\) 的像，是紧凸集；它包含零函数和 \(\widehat{\mathrm{G}}\) 中的元素（事实上，由 II, p. 220, 定理 2，这些元素都具有 \(\mathrm{ev}_x\colon \chi \mapsto \chi (x)\) 的形式，其中 \(x\) 是 G 中某个元，且 \(\mathrm{ev}_x = \mathcal{F}(\varepsilon_{x^{-1}})\)）。因此集合 \(\mathcal{F}(\mathcal{N})\) 包含 \(\mathrm{Pos}_{\leqslant 1}(\mathrm{G})\) 的极点，由克赖因–米尔曼定理（EVT, II, p. 59, 定理 1）可得 \(\mathcal{F}(\mathcal{N}) = \mathrm{Pos}_{\leqslant 1}(\mathrm{G})\)。证明完毕。

当 \(G = R^{k}\)，其中整数 \(k \in N\) 时，本定理对应于 INT, IX, p. 94, § 6, no. 12, 命题 11。

